%%%%%%%%%%%%%%%%%%%%%%%%%%%%%%%%%%%%%%%%%
% Medium Length Professional CV
% LaTeX Template
% Version 2.0 (8/5/13)
%
% This template requires the resume.cls file to be in the same directory as the
% .tex file. The resume.cls file provides the resume style used for structuring the
% document.
%
%%%%%%%%%%%%%%%%%%%%%%%%%%%%%%%%%%%%%%%%%

%----------------------------------------------------------------------------------------
%	PACKAGES AND OTHER DOCUMENT CONFIGURATIONS
%----------------------------------------------------------------------------------------

\documentclass{resume} % Use the custom resume.cls style
\usepackage[colorlinks=true,
            linkcolor=red,
            urlcolor=blue,
            citecolor=gray]{hyperref}

\usepackage[left=0.75in,top=0.75in,right=0.75in,bottom=0.75in]{geometry} % Document margins
\newcommand{\tab}[1]{\hspace{.2667\textwidth}\rlap{#1}}
\newcommand{\itab}[1]{\hspace{0em}\rlap{#1}}
\name{Yitian Wang} % Your name
%\address{\href{http://linkedin.com/in/tim-wang-yitian/}{http://linkedin.com/in/tim-wang-yitian/}} % Your address
%\address{123 Pleasant Lane \\ City, State 12345} % Your secondary addess (optional)
%\address{\href{mailto:ywang1057@ucr.edu}{ywang1057@ucr.edu} \\ +1(929)624-1021 \\ \href{http://cosmotim.github.io}{Portfolio Website} \\ \href{http://linkedin.com/in/tim-wang-yitian/}{Linkedin}} % Your phone number and email
\begin{document}

% Two-column contact layout
\vspace{-1em}
\noindent
\begin{minipage}[t]{0.5\textwidth}
\raggedright
+1(929)624-1021 \\
\href{mailto:ywang1057@ucr.edu}{ywang1057@ucr.edu}
\end{minipage}%
\begin{minipage}[t]{0.5\textwidth}
\raggedleft
\href{http://cosmotim.github.io}{cosmotim.github.io} \\
\href{http://linkedin.com/in/tim-wang-yitian/}{linkedin.com/in/tim-wang-yitian}
\end{minipage}
\vspace{0.5em}

%----------------------------------------------------------------------------------------
%	SUMMARY SECTION
%----------------------------------------------------------------------------------------

\begin{rSection}{SUMMARY}
Ph.D. scholar in Electrical Engineering specializing in thermal transport and lattice dynamics in functional oxides and solid-state ionic conductors.
Experienced in crystal and thin-film growth, vacuum deposition, material characterization, and data-driven modeling using Python and MATLAB.
Published 5 first-author works in leading journals and skilled at collaborating with national labs and cross-disciplinary device and materials teams.
Additional teaching assistant experience during Ph.D. studies.
\end{rSection}

%----------------------------------------------------------------------------------------
%	EDUCATION SECTION
%----------------------------------------------------------------------------------------

\begin{rSection}{EDUCATION}
{\bf University of California, Riverside, CA}\hfill{Mar.\ 2021 -- Jun.\ 2026 (expected)}
\\ Ph.D.\ in Electrical and Computer Engineering
\\
\\{\bf Columbia University in the City of New York, NY}\hfill{Aug.\ 2019 -- Feb.\ 2021}
\\ Master of Science, Material Science and Engineering
\\
\\{\bf University of California, Berkeley, CA}\hfill{Jul.\ 2017 -- Aug.\ 2017}
\\ Summer Session, Physics
\\
\\{\bf Beijing Normal University, CN}\hfill{Sep.\ 2015 -- Jun.\ 2019}
\\ Bachelor of Science, Physics

\end{rSection}


\begin{rSection}{SKILLS}

\begin{tabular}{ @{} >{\bfseries}l @{\hspace{6ex}} p{12.5cm} }
Software & Python, MATLAB, R, SQL, LaTeX, COMSOL, SolidWorks, Linux \\[4pt]
Analysis Tools & GSAS-II, Mantid, ImageJ, VESTA, Quantum ESPRESSO, VASP \\[4pt]
Synthesis \& Growth & Floating-Zone Crystal Growth, PVD / Magnetron Sputtering, CVD \\[4pt]
Characterization & SEM, TEM, Scanning Tunneling Microscopy (STM), XRD, AFM, Inelastic Neutron Scattering, Raman, FTIR \\[4pt]
Transport Properties & PPMS, DSC, TGA, Four-Probe Transport, Mechanical Strength Testing \\[4pt]
Battery \& Electrochem & Prototype Cell Assembly, Glovebox Operation, Separator Fabrication, EIS, Galvanostatic Cycling, Thermal Abuse Testing \\[4pt]
Methods & Design of Experiments (DOE), AI Agent Training
\end{tabular}

\end{rSection}

\begin{rSection}{AWARDS \& ACTIVITIES}
\begin{itemize}
\itemsep 0.1em
\item Dissertation Completion Fellowship Award, University of California, Riverside (2025)
\item MRS 2023 Spring Meeting Highly Commended Student Talk Award
\item `Jingshi' Scholarship, Beijing Normal University (2016 -- 2018)
\item COMAP Mathematical Contest in Modeling --- team member and team leader
\end{itemize}
\end{rSection}

\begin{rSection}{PROFESSIONAL MEMBERSHIPS}
\begin{itemize}
\itemsep 0.1em
\item Sigma Xi, The Scientific Research Honor Society --- Elected Full Member (2026)
\item American Physical Society (APS) --- Member (2025 -- 2026)
\item Materials Research Society (MRS) --- Student Member (through 2027)
\end{itemize}
\end{rSection}

%--------------------------------------------------------------------------------
%    Research Experience
%-----------------------------------------------------------------------------------------------

\begin{rSection}{RESEARCH EXPERIENCE}

\begin{rSubsection}
{Graduate Research Assistant}{Jun.\ 2021 -- Jun.\ 2025}
{University of California}{Riverside, CA}
\item Pioneered lattice-dynamics studies on Li-ion solid-state electrolytes, extending work previously limited to Cu- and Ag-ion systems.
\item Published 5 first-author papers in journals including \textit{PRX Energy} and \textit{J.\ Mater.\ Chem.\ A}.
\item Built a floating-zone crystal growth workflow from scratch, producing centimeter-scale LLZTO single crystals that enabled phonon-resolved neutron scattering measurements previously infeasible on polycrystalline samples.
\item Identified \textit{diffuson}-mediated thermal transport in Li-ion conductors, providing the first verification of this mechanism in Li-ion solid-state electrolytes.
\item Developed a two-channel (phonon + \textit{diffuson}) thermal-conductivity model in MATLAB that explained anomalous temperature dependence in Li-ion conductors where the traditional Debye model failed.
\end{rSubsection}

\begin{rSubsection}
{Experiment Lead}{Jan.\ 2022 -- Jul.\ 2023}
{Oak Ridge National Laboratory}{Oak Ridge, TN}
\item Authored two successful beam time proposals and led inelastic neutron-scattering experiments on LLZTO single crystals using TAX (HFIR) and ARCS (SNS) at Oak Ridge National Laboratory (ORNL).
\item Resolved the full phonon dispersion of a garnet solid electrolyte for the first time, providing experimental benchmarks for ab initio lattice-dynamics calculations.
\item Built automated data-reduction pipelines in Mantid and Python, significantly reducing manual analysis time.
\item Managed on-site experimental decisions and coordinated across $\sim$7 research groups spanning UCR, UT Austin, ORNL, ETH Z\"{u}rich, and TU Wien.
\end{rSubsection}

\begin{rSubsection}
{Student Researcher}{Sep.\ 2019 -- Jan.\ 2021}
{Columbia University}{New York, NY}
\item Engineered composite Li-ion battery separators ($\gamma$-C$_3$N$_4$/PVDF) via solution processing, targeting improved electrochemical reaction kinetics.
\item Assembled and tested CR2032 full cells in an Ar-filled glovebox over 100+ galvanostatic cycles, characterizing cycling performance and interfacial stability via EIS.
\item Applied first-principles DFT (Quantum ESPRESSO) to model defect stability in olivine cathode structures, complementing experimental work with computational materials design.
\end{rSubsection}

\begin{rSubsection}
{Student Researcher}{Sep.\ 2018 -- May 2019}
{Beijing Normal University}{Beijing, CN}
\item Designed a uniaxial-strain apparatus in SolidWorks exploiting differential thermal contraction for controlled detwinning of iron-based superconductors below 140 K.
\item Built a Monte Carlo simulation in MATLAB to model atomic movement of metal atoms inside multi-wall carbon nanotubes, reproducing experimental observations.
\end{rSubsection}

\end{rSection}

%--------------------------------------------------------------------------------
%    Teaching Experience
%-----------------------------------------------------------------------------------------------

\begin{rSection}{TEACHING EXPERIENCE}

{\bf Graduate Teaching Assistant} \hfill {\bf University of California, Riverside}

\vspace{0.3em}
\begin{tabular}{ @{} p{8.5cm} l r l }
Technology in the Premodern World (ENGR 108) & Reader & Fall 2025 & \textit{N/A} \\
Circuits and Electronics Lab (EE 005) & Lab Instructor & Spring 2025 & \textit{7.0/7.0 (96th \%ile)} \\
Technology in the Premodern World (ENGR 108) & TA / Reader & Spring 2025 & \textit{5.7/7.0 (83rd \%ile)} \\
Linear Methods for Engr.\ Analysis -- MATLAB (EE 020B) & Discussion TA & Winter 2023 & \textit{6.5/7.0 (81st \%ile)} \\
Magnetic Materials (EE 139) & Discussion TA & Fall 2022 & \textit{6.5/7.0 (77th \%ile)} \\
\end{tabular}
\vspace{0.3em}

\begin{list}{$\cdot$}{\leftmargin=0em}
\itemsep -0.5em \vspace{-0.5em}
\item Served 280+ students across 7 sections over 3 academic years.
\item Achieved perfect 7.0/7.0 evaluation (96th percentile campus-wide) in EE 005, with consistently above-department averages across all courses.
\item Selected student feedback: ``singlehandedly one of the most effective teachers on this campus'' and ``top tier educator'' (EE 005); ``the backbone of the entire class'' (EE 020B); ``always helpful... feedback was helpful in creating a well organized presentation'' (EE 139).
\end{list}
\vspace{0.5em}

\end{rSection}

%--------------------------------------------------------------------------------
%    Internship Experience
%-----------------------------------------------------------------------------------------------

\begin{rSection}{INTERNSHIP EXPERIENCE}

\begin{rSubsection}
{Administrative Assistant, Physics Department}{Sep.\ 2018 -- Jan.\ 2019}
{Beijing Normal University}{Beijing, CN}
\item Managed archival of 3 years of departmental paperwork and organized the academic calendar for 4 course sections.
\end{rSubsection}

\begin{rSubsection}
{Student Intern}{Jul.\ 2016 -- Aug.\ 2016}
{Shanghai East China Computer Corp.}{Shanghai, CN}
\item Gained hands-on exposure to IT industry workflows, strengthening teamwork and cross-functional communication skills.
\end{rSubsection}

\end{rSection}

%	Publications
%----------------------------------------------------------------------------------------
\begin{rSection}{SELECTED PUBLICATIONS}
\begin{enumerate}
\itemsep 0.2em

\item \textbf{Y. Wang}, Q. Jia, S. Li, L. Shi, Y. Li, X. Chen.
``Low thermal conductivity and lattice anharmonicity of NaSICON-type solid electrolyte Na$_3$Zr$_2$Si$_2$PO$_{12}$.''
\textit{Tungsten}, Nov.\ 2025.
\href{https://doi.org/10.1007/s42864-025-00357-6}{[Link]}

\item \textbf{Y. Wang}, Q. Jia, S. Li, L. Shi, Y. Li, X. Chen.
``Glass-like thermal transport in polycrystalline perovskite Li-ion conductor Li$_{3/8}$Sr$_{7/16}$Hf$_{1/4}$Ta$_{3/4}$O$_3$.''
\textit{Chem.\ Commun.}, Sep.\ 2025.
\href{https://doi.org/10.1039/D5CC04693A}{[Link]}
\href{https://doi.org/10.1039/D5CC90366A}{[Front Cover]}

\item \textbf{Y. Wang}, Y. Su, J. Carrete, H. Zhang, N. Wu, \textit{et al.}
``Origin of intrinsically low thermal conductivity in a garnet-type solid electrolyte.''
\textit{PRX Energy}, Jul.\ 2025 (OA fee waived).
\href{https://doi.org/10.1103/6wj2-kzhh}{[Link]}

\item S. Li, S. Guo, \textbf{Y. Wang}, H. Li, Y. Xu, V. Carta, J. Zhou, X. Chen.
``Enhanced magnon thermal transport in yttrium-doped spin ladder compounds Sr$_{14-x}$Y$_x$Cu$_{24}$O$_{41}$.''
\textit{J.\ Appl.\ Phys.}, Jul.\ 2024.
\href{https://doi.org/10.1063/5.0214897}{[Link]}

\item \textbf{Y. Wang}, S. Li, N. Wu, Q. Jia, T. Hoke, L. Shi, Y. Li, X. Chen.
``Thermal properties and lattice anharmonicity of Li-ion conducting garnet solid electrolyte Li$_{6.5}$La$_3$Zr$_{1.5}$Ta$_{0.5}$O$_{12}$.''
\textit{J.\ Mater.\ Chem.\ A}, Jun.\ 2024.
\href{https://doi.org/10.1039/D4TA02264E}{[Link]}

\item S. Guo, H. Li, X. Bai, \textbf{Y. Wang}, S. Li, R. E. Dunin-Borkowski, J. Zhou, X. Chen.
``Size-dependent magnon thermal transport in a nanostructured quantum magnet.''
\textit{Cell Rep.\ Phys.\ Sci.}, Mar.\ 2024.
\href{https://doi.org/10.1016/j.xcrp.2024.101879}{[Link]}

\item Y. Xu, Z. Barani, P. Xiao, S. Sudhindra, \textbf{Y. Wang}, \textit{et al.}
``Crystal structure and thermoelectric properties of layered van der Waals semimetal ZrTiSe$_4$.''
\textit{Chem.\ Mater.}, Sep.\ 2022.
\href{https://doi.org/10.1021/acs.chemmater.2c02155}{[Link]}

\item \textbf{Y. Wang}, X. Chen.
``Single crystal growth and electrochemical studies of garnet-type fast Li-ion conductors.''
\textit{Tungsten}, Sep.\ 2022.
\href{https://doi.org/10.1007/s42864-022-00176-z}{[Link]}

\end{enumerate}
\end{rSection}

%----------------------------------------------------------------------------------------
%	MEDIA COVERAGE
%----------------------------------------------------------------------------------------
\begin{rSection}{MEDIA COVERAGE}
\begin{itemize}
\itemsep 0.1em
\item \href{https://www.youtube.com/watch?v=dPwabntYj9s}{Why LLZTO is the Secret to Super-Fast EV Charging} --- UC Riverside, Feb.\ 2026
\item \href{https://news.ucr.edu/articles/2025/10/22/cool-battery-power}{Cool Battery Power} --- UC Riverside News, Oct.\ 2025
\item \href{https://techxplore.com/news/2025-10-solid-electrolyte-unique-atomic-generation.html}{Solid electrolyte's unique atomic structure helps next-generation batteries keep their cool} --- Tech Xplore, Oct.\ 2025
\item \href{https://ece.ucr.edu/news/2023/05/10/ece-phd-student-wins-materials-research-student-talk-award}{ECE PhD Student Wins Materials Research Student Talk Award} --- UCR ECE Department, May 2023
\end{itemize}
\end{rSection}

%----------------------------------------------------------------------------------------

 \end{document}
